\section{Preliminaries}
\label{sec:prelim}

Throughout, \(\mathbb{F}\) is a prime field and \(\lambda\) is the security
parameter.  Matrices are denoted by bold uppercase letters and vectors by bold
lowercase letters.  For \(\mat{M}\), we write \(\mat{M}[i,*]\) and
\(\mat{M}[*,j]\) for its \(i\)-th row and \(j\)-th column.  We use
\([n]=\{0,\ldots,n-1\}\), \(I\sample[n]^t\) for \(t\) uniform query
positions, and \(r\sample\mathbb{F}\) for a uniform field element.

\subsection{Linear Codes and Folding}
\label{sec:prelim-codes}

Let \(\mathcal{C}:\mathbb{F}^k\to\mathbb{F}^n\) be a linear code with
generator matrix \(\mat{G}\), rate \(\rho=k/n\), and relative distance
\(\delta\).  We write \(\enc(\vect{m})=\vect{m}\mat{G}\).  A matrix
\(\mat{M}\in\mathbb{F}^{k\times k}\) is encoded row-wise as
\(\widehat{\mat{M}}=\enc(\mat{M})\in\mathbb{F}^{k\times n}\).

\begin{definition}[Interleaved code]\label{def:interleaved-code}
The \(\ell\)-fold interleaved code \(\mathcal{C}^{\langle\ell\rangle}\)
groups \(\ell\) codewords position by position.  Its \(j\)-th symbol is an
element of \(\mathbb{F}^{\ell}\).  In particular, a row-wise encoded
\(k\times k\) matrix is a \(k\)-fold interleaved codeword whose \(j\)-th
symbol is \(\widehat{\mat{M}}[*,j]\).
\end{definition}

For coefficients \(r_0,\ldots,r_{k-1}\in\mathbb{F}\), code linearity gives
\begin{equation}\label{eq:code-linearity}
  \sum_{i=0}^{k-1}r_i\widehat{\mat{M}}[i,*]
  =\enc\!\left(\sum_{i=0}^{k-1}r_i\mat{M}[i,*]\right).
\end{equation}
For \(\vect{r},\vect{c}\in\mathbb{F}^k\), define
\begin{equation}\label{def:fold}
  \fold(\vect{r},\vect{c})\coloneqq
  \langle\vect{r},\vect{c}\rangle.
\end{equation}
Thus \(\fold(\vect{r},\widehat{\mat{M}}[*,j])\) is the \(j\)-th coordinate
of \(\enc(\vect{r}\mat{M})\).  Our implementation uses a Reed--Solomon code;
Appendix~\ref{app:rs-barycentric} gives its concrete consistency check.

\subsection{Commitment and Proof Interfaces}
\label{sec:prelim-commitments}

We use Merkle trees as position-binding vector commitments and Pedersen
commitments as hiding, computationally binding, and additively homomorphic
commitments.  We write \(\merkleC\), \(\merkleO\), and \(\merkleV\) for
Merkle commitment, opening, and verification, and \(\pedC\) for Pedersen
commitment.  Merkle leaves in our construction are Pedersen commitments, so
the tree itself uses deterministic leaves.

\label{sec:prelim-snark}
A commit-and-prove SNARK (CP-SNARK)~\cite{gnark,legosnark} proves an NP
relation while binding a designated witness part \(\comwit\) to a public
witness commitment \(\widetilde{\cm}=\bPi_{\cp}.\com(\para,\comwit;
\widetilde{o})\).  We use the interface
\(\bPi_{\cp}=(\setup,\com,\prove,\verify)\).

\label{sec:prelim-cplink}
CP-Link proves that a SNARK-side commitment to an ordered concatenation of
sampled blocks contains the same messages as the corresponding external
Pedersen commitments.  For blocks \((\vect{m}_i)_{i=0}^{q-1}\), it links
\[
\begin{aligned}
  \widetilde{\cm}
  &=\com(\ck_S,\vect{m}_0\|\cdots\|\vect{m}_{q-1};\widetilde{o}),\\
  \cm_i&=\com(\ck_E,\vect{m}_i;o_i)\quad(i\in[q]).
\end{aligned}
\]
The construction requires CP-SNARK soundness and zero knowledge, commitment
binding and hiding, and CP-Link soundness and witness zero knowledge; these
assumptions are summarized in Definition~\ref{def:lamp-backend-assumptions}.

\subsection{Structured Freivalds Check}
\label{sec:prelim-freivalds}

Freivalds' algorithm~\cite{freivalds1977probabilistic} verifies
\(\mat{A}\mat{B}=\mat{C}\) by checking
\(\vect{r}\mat{A}\mat{B}=\vect{r}\mat{C}\).  We use the compact structured
challenge \(\vect{r}=(1,r,\ldots,r^{k-1})\) for
\(r\sample\mathbb{F}\).  If
\(\mat{D}=\mat{A}\mat{B}-\mat{C}\neq0\), a nonzero column of \(\mat{D}\)
defines a nonzero univariate polynomial in \(r\) of degree at most \(k-1\).
By Schwartz--Zippel~\cite{schwartz1980fast,zippel1979probabilistic},
\[
  \Pr[\vect{r}\mat{D}=0]\le\frac{k-1}{|\mathbb{F}|}.
\]
